% =========================================================
% Compléments M1 — Lemme d'Itô (C. Dérivations)
% =========================================================

\subsection*{(1) Preuve guidée (dimension 1) : où apparaît exactement le terme $\tfrac12 f_{xx}\sigma^2$}
On donne ici une preuve structurée (niveau M1) de la formule d'Itô en dimension 1, en suivant l'idée «partition + Taylor + variation quadratique».
Soit $X$ un processus d'Itô :
\[
X_t=X_0+\int_0^t a_s\,\dd s+\int_0^t b_s\,\dd W_s.
\]
Fixons une partition $\pi:0=t_0<\dots<t_n=t$ et écrivons, par télescopage,
\[
f(t,X_t)-f(0,X_0)=\sum_{k=0}^{n-1}\Big(f(t_{k+1},X_{t_{k+1}})-f(t_k,X_{t_k})\Big).
\]
On ajoute et retranche $f(t_{k+1},X_{t_k})$ :
\[
\Delta f_k
=\underbrace{f(t_{k+1},X_{t_k})-f(t_k,X_{t_k})}_{(\mathrm{I})}
\;+\;\underbrace{f(t_{k+1},X_{t_{k+1}})-f(t_{k+1},X_{t_k})}_{(\mathrm{II})}.
\]

\paragraph{(I) Variation en temps.}
Par Taylor en temps :
\[
(\mathrm{I})=\partial_t f(t_k,X_{t_k})\Delta t_k + r_k^{(t)},
\]
où $\sum_k r_k^{(t)}\to 0$ quand $|\pi|\to 0$ (régularité $C^{1,2}$).

\paragraph{(II) Variation en espace.}
Par Taylor en espace à l'ordre 2 :
\[
(\mathrm{II})
=\partial_x f(t_{k+1},X_{t_k})\Delta X_k
\frac12\partial_{xx}f(t_{k+1},X_{t_k})(\Delta X_k)^2+r_k^{(x)}.
\]

\paragraph{Contrôle des ordres de grandeur.}
On a
\[
\Delta X_k = \int_{t_k}^{t_{k+1}} a_s\,\dd s + \int_{t_k}^{t_{k+1}} b_s\,\dd W_s.
\]
Le premier terme est $O(\Delta t_k)$, le second est $O(\sqrt{\Delta t_k})$.
Donc $(\Delta X_k)^2$ est dominé par la partie brownienne et contribue à l'ordre $\Delta t_k$ (c'est le c\oe{}ur du raisonnement).

\paragraph{Passage à la limite : identification des trois termes.}
En sommant :
\begin{itemize}[leftmargin=*]
\item Les sommes des termes $\partial_t f\,\Delta t_k$ convergent vers $\int_0^t \partial_t f(s,X_s)\,\dd s$.
\item Les sommes des termes $\partial_x f\,\Delta X_k$ convergent vers $\int_0^t \partial_x f(s,X_s)\,\dd X_s$ (qui se décompose ensuite en drift + intégrale d'Itô).
\item La somme des termes $\frac12\partial_{xx}f\,(\Delta X_k)^2$ converge vers $\frac12\int_0^t \partial_{xx} f(s,X_s)\,\dd[X]_s$.
\end{itemize}
Le point crucial est le troisième item : comme $[X]_s=\int_0^s b_u^2\,\dd u$, on obtient le terme
\[
\frac12\int_0^t \partial_{xx}f(s,X_s)b_s^2\,\dd s,
\]
qui est \emph{absent} en calcul classique.
Les restes $\sum_k r_k^{(t)}$ et $\sum_k r_k^{(x)}$ tendent vers 0 par régularité.
On obtient la formule d'Itô.

\subsection*{(2) Exemple canonique : martingale exponentielle via Itô}
Reprenons
\[
Z_t=\exp\Big(\theta W_t-\tfrac12\theta^2 t\Big).
\]
Posons $f(t,w)=\exp(\theta w-\tfrac12\theta^2 t)$.
Alors
\[
\partial_t f=-\tfrac12\theta^2 f,\qquad \partial_w f=\theta f,\qquad \partial_{ww}f=\theta^2 f.
\]
Par Itô :
\[
\dd Z_t = \Big(\partial_t f+\tfrac12\partial_{ww}f\Big)\dd t + \partial_w f\,\dd W_t
=\Big(-\tfrac12\theta^2 f+\tfrac12\theta^2 f\Big)\dd t + \theta f\,\dd W_t
=\theta Z_t\,\dd W_t.
\]
Le drift est nul : $Z$ est une martingale. C'est exactement la forme utilisée dans Girsanov.

\subsection*{(3) Exemple finance : actualisation et produit stochastique}
Soit $B_t=e^{rt}$ le compte monétaire.
La variable actualisée $\widetilde S_t=e^{-(r-q)t}S_t$ doit être une martingale sous $\Qmeas$ en BS.
Vérifions-le par calcul :
\[
\dd S_t=(r-q)S_t\,\dd t+\sigma S_t\,\dd W_t^{\Qmeas},
\qquad
\dd\big(e^{-(r-q)t}\big)=-(r-q)e^{-(r-q)t}\,\dd t.
\]
En appliquant la formule produit $\,\dd(UV)=U\,\dd V+V\,\dd U+\dd[U,V]\,$ avec $U=e^{-(r-q)t}$ (variation finie, donc crochet nul), on obtient
\[
\dd\widetilde S_t
=e^{-(r-q)t}\dd S_t+S_t\,\dd(e^{-(r-q)t})
=e^{-(r-q)t}\sigma S_t\,\dd W_t^{\Qmeas},
\]
donc $\widetilde S$ est bien une martingale.

\subsection*{(4) Exercices corrigés (Itô)}
\paragraph{Exercice 1 --- SDE de $1/S_t$ sous GBM.}
Sous $\dd S_t=(r-q)S_t\,\dd t+\sigma S_t\,\dd W_t$, calculer la dynamique de $Y_t=1/S_t$.

\emph{Solution.}
Poser $f(s)=s^{-1}$, alors $f'(s)=-s^{-2}$ et $f''(s)=2s^{-3}$.
Itô :
\[
\dd Y_t=f'(S_t)\dd S_t+\tfrac12 f''(S_t)(\dd S_t)^2.
\]
Or $(\dd S_t)^2=\sigma^2 S_t^2\,\dd t$.
Donc
\[
\dd Y_t=-\frac{1}{S_t^2}\big((r-q)S_t\,\dd t+\sigma S_t\,\dd W_t\big)
\frac12\cdot 2\frac{1}{S_t^3}\sigma^2 S_t^2\,\dd t
=-(r-q)Y_t\,\dd t-\sigma Y_t\,\dd W_t+\sigma^2 Y_t\,\dd t,
\]
soit
\[
\boxed{\dd Y_t=(\sigma^2-(r-q))Y_t\,\dd t-\sigma Y_t\,\dd W_t.}
\]

\paragraph{Exercice 2 --- calcul de $\E[S_T^p]$ (rappel).}
Montrer à partir de la dynamique de $S_t^p$ que
\[
\E^{\Qmeas}[S_T^p]=S_0^p\exp\Big(p(r-q)T+\tfrac12 p(p-1)\sigma^2T\Big).
\]
\emph{Solution.}
On a déjà obtenu
\[
\dd(S_t^p)=S_t^p\Big(p(r-q)+\tfrac12 p(p-1)\sigma^2\Big)\dd t + p\sigma S_t^p\,\dd W_t.
\]
En prenant l'espérance, le terme martingale disparaît, donc
\[
\frac{\dd}{\dd t}\E[S_t^p]=\Big(p(r-q)+\tfrac12 p(p-1)\sigma^2\Big)\E[S_t^p],
\]
ODE résolue par exponentielle, avec condition initiale $\E[S_0^p]=S_0^p$.
